% =============================================================================
% Seção 5 — Discussão
% =============================================================================

\section{Discussão}
\label{sec:discussao}

\subsection{Acessibilidade de Hardware e Contexto de Implantação}

Uma restrição central de design do NeuroAssist AI foi a implantabilidade em hardware
comum já presente na maioria dos domicílios. O pipeline de detecção requer apenas
uma webcam padrão (720p ou superior) e um navegador moderno; não são necessárias GPU,
hardware de sensoriamento dedicado ou dongle proprietário. Isso contrasta com os
dispositivos comerciais de CAA por rastreamento ocular, cujo custo tipicamente excede
USD~3.000~\cite{beukelman2013aac} e cujo setup requer envolvimento de técnico clínico.
Uma webcam adequada para o pipeline custa aproximadamente BRL~100--500 e está
amplamente disponível em países de baixa e média renda (LMICs), incluindo o Brasil,
onde aproximadamente 30.000 indivíduos vivem com ELA~\cite{mehta2018mnd}. O design
offline-first elimina adicionalmente a dependência de banda larga estável — uma
barreira significativa em contextos rurais e periurbanos em LMICs — uma vez que
nenhuma conectividade à internet é necessária durante a detecção de gestos ou a
entrega de alarmes em uma rede local.

\subsection{Privacidade desde a Concepção}

O vídeo facial é classificado como dado biométrico tanto pelo Regulamento Geral de
Proteção de Dados (GDPR) quanto pela Lei Geral de Proteção de Dados Pessoais (LGPD),
atraindo o mais alto nível de proteção regulatória para contextos de
saúde~\cite{edpb2020privacy}. O NeuroAssist AI satisfaz um princípio estrito de
privacidade desde a concepção: o runtime WebAssembly e o modelo de landmarks faciais
são executados inteiramente dentro do ambiente de execução em sandbox do navegador, e
nenhum quadro de vídeo ou coordenada de landmark facial é transmitida pela rede. O
único dado que cruza a fronteira de rede é um payload JSON estruturado contendo o tipo
de gesto detectado, uma pontuação de confiança e um identificador de paciente — nenhum
deles constitui dado biométrico segundo as definições regulatórias atuais. Essa
arquitetura torna o NeuroAssist AI adequado para implantação domiciliar direta sem
exigir acordos institucionais de processamento de dados.

\subsection{Extensibilidade via Mecanismo Plug-in \texttt{engineFeatureRef}}

O campo \texttt{engineFeatureRef} no \texttt{DetectionProfile} foi projetado
explicitamente para expansão futura de gestos. Adicionar um novo tipo de gesto —
como elevação de canto labial, fechamento ocular sustentado, ou aceno de cabeça via
estimativa de pose — requer apenas: (i) inserir uma nova linha \texttt{TipoGesto}
com um novo valor \texttt{engineFeatureRef}; (ii) implementar o módulo detector
correspondente no agente de borda; e (iii) publicar um novo \texttt{DetectionProfile}.
Nenhum modelo de domínio, caso de uso ou regra ArchUnit existente requer modificação.

\subsection{Limitações}

\textbf{Dependência de iluminação.} O EAR é calculado a partir de coordenadas de
marcos no espaço de pixels; sua confiabilidade se deteriora em ambientes com pouca
iluminação. Testes práticos documentaram casos onde o $\theta_\text{close}$ padrão de
0,20 exigiu ajuste manual para 0,17--0,18 sob iluminação artificial fraca. Trabalhos
futuros devem incorporar ajuste automático de limiar adaptativo à iluminação.

\textbf{Escopo do estudo de caso preliminar.} A avaliação inclui um paciente com ELA
em estágio avançado e cinco cuidadores e voluntários, e deve ser interpretada como
estudo de caso preliminar, e não como ensaio clínico. A generalização para a população
mais ampla de ELA exigirá um estudo multicêntrico com amostra estatisticamente
representativa sob aprovação de CEP.

\textbf{Sobrecarga de latência por polling.} A aplicação do cuidador detecta eventos
ativos consultando o \texttt{evnt-service} a cada cinco segundos. Isso introduz uma
sobrecarga máxima fixa de 5\,s na latência evento-para-alarme visível. O intervalo de
polling é uma simplificação do PoC; a implantação em produção deve substituir o polling
por eventos enviados pelo servidor (SSE) ou push WebSocket.

\textbf{Adaptadores stub para módulos periféricos.} O catálogo de comandos (serviço
PAC), a trilha de auditoria (serviço HIST) e o cadastro de cuidadores (serviço ACS)
são representados por adaptadores stub com valores fixos no PoC. Os resultados de
latência e confiabilidade refletem apenas os subsistemas GEST e EVNT.

\subsection{Trabalhos Futuros}

Os próximos passos mais impactantes, em ordem de prioridade clínica, são: (i)~um
estudo clínico com participantes com ELA sob aprovação de CEP; (ii)~substituição do
transporte de polling por SSE para eliminar o teto de latência de alarme de 5\,s;
(iii)~integração do catálogo completo de comandos PAC para suportar comunicação
multi-comandos; e (iv)~avaliação de tipos de gesto adicionais, particularmente para
pacientes que retêm controle unilateral de sobrancelha após a perda do controle
bilateral de piscada.

Um caminho de integração de alta prioridade é o acoplamento dos eventos de angústia
do NeuroAssist AI com fluxos de dados de ventiladores VNI. Grupos de pesquisa que
estudam PVA na ELA~\cite{stell2024pva, yang2026pva} trabalham com sinais fisiológicos
multicanais em formato EDF, com precisão temporal de segundos. O \texttt{eventoId}
do NeuroAssist AI carrega um timestamp \texttt{createdAt} do servidor e um timestamp
\texttt{recognizedAt} do dispositivo; correlacionar esses dados com a linha temporal
de monitoramento da VNI permitiria a um estudo futuro testar se eventos de angústia
confirmados pelo cuidador co-ocorrem com índices elevados de PVA — uma hipótese
clinicamente significativa que nenhum dos dois sistemas pode testar isoladamente.
